\documentclass[a4paper,11pt]{article} % Don't change this
\usepackage{fullpage} % Don't change this
\usepackage{titling} % Don't change this
\pretitle{\vspace{-6em}\begin{center}\larger[2]\bf} % Don't change this
\postauthor{\end{center}} % Don't change this
\preauthor{\vspace{-2em}\begin{center}\smaller[1]} % Don't change this
\postauthor{\end{center}\vspace{-5em}} % Don't change this
\usepackage[shortlabels]{enumitem} % Don't change this
\setlist{itemsep=0em,topsep=0.3em} % Don't change this

%%% Some useful packages.
\usepackage{xcolor}
\usepackage{float}
\usepackage{amsmath,amssymb,mathtools}
\usepackage{graphicx}
\usepackage{subcaption}
\usepackage[hidelinks]{hyperref}
\usepackage{relsize}
\usepackage{cleveref}
\crefname{figure}{Fig.}{Figs.}
\usepackage[backend=bibtex,natbib=true,style=authoryear-icomp,maxnames=3,giveninits=true]{biblatex}
\citetrackerfalse
\bibliography{ref} % put bib entries in ref.bib
\usepackage[linesnumbered]{algorithm2e}
\usepackage{cleveref}
\crefname{equation}{Equation.}{Equations.}
\crefname{figure}{Figure}{Figures.}
\usepackage{lipsum}
\usepackage{wrapfig}
\usepackage{tikz}

%%% Add additional packages or macros below

\begin{document}

\title{Forecast-Informed Deep Learning for Electricity Price Prediction in the Australian NEM}

\author{Group 12 \\\\
Connor Kikkert (s4748070) $\cdot$
Connor McAreavey (s4884456) $\cdot$
Tyreece Paul (s4883782) \\\\
John Pioc (s4775725) $\cdot$ Massimo Spadaro (s4850133)}

\date{}
\maketitle

\begin{abstract}
Spot electricity prices in Australia's National Electricity Market (NEM) are highly
volatile, non-stationary, and prone to extreme price spikes. Using publicly available
AEMO market and forecast data, we propose to implement and evaluate deep learning
architectures for price forecasting across the five NEM regions, building from the
recurrent and convolutional models covered in this course towards a graph-attention,
multi-task architecture inspired by \citet{li2025mgaal}, benchmarked against the
comparison of \citet{gani2026deeptime}.
\end{abstract}

\section{Topic}

We propose to implement a deep-learning based short-term electricity price forecasting
system for the Australian NEM, which is composed of five interconnected regional markets
(NSW, QLD, SA, TAS, VIC). We will adapt a simplified version of the multi-task graph
adaptive learning (MGAAL) model of \citet{li2025mgaal}, which frames price forecasting
as a multi-task problem: a primary regression task predicting a future regional price,
and an auxiliary classification task predicting anomalous price-spike events, with the
two tasks sharing graph-attention and LSTM-based spatiotemporal feature extraction. The
graph-attention component treats each NEM region's price series as a node in a graph and
learns attention weights that capture relationships between regional markets, while LSTM
layers capture longer-range temporal dependencies. Historical realised market data will
be obtained using \texttt{NEMOSIS}, while historical AEMO ahead-process forecast data
will be obtained using \texttt{NEMSEER} and aligned by forecast issue time to avoid
look-ahead leakage. We will then compare the designed model against a set of standard
deep learning baselines (MLP, CNN, LSTM, CNN-LSTM, and a self-attention/Transformer
model) trained on the same data splits, forecast horizons (day-ahead and two-day-ahead),
and evaluation metrics (MAE, RMSE, sMAPE, and directional accuracy), with particular
attention to performance during extreme and negative price events. This methodology is
inspired by the benchmarking approach of \citet{gani2026deeptime}, and directly tests
whether an increase in model complexity correlates with better performance on volatile,
regime-switching electricity price series, while relating our findings to concepts
covered in the course.

\section{Significance}

Electricity is sold at unregulated spot prices but bought by consumers at regulated
rates, so when prices swing wildly, retailers take the hit directly, which can threaten3 Fea
both their finances and the stability of the grid \citep{li2025mgaal}. Between October
2021 and January 2022, average NEM spot prices were almost four times higher than the
year before and peaked at \$14{,}000/MWh \citep{li2025mgaal}. More recent data from
2023--2025 shows this hasn't gone away: increased renewable penetration has contributed
to negative prices, steeper ramps, and higher short-term volatility, with intraday
analysis showing elevated midday negative-price occurrence and evening volatility as
solar output declines \citep{gani2026deeptime}. Being able to forecast prices accurately
in these conditions is useful for retailers managing risk, generators deciding when to
bid and system operators trying to keep enough reserve capacity on hand.
\citet{gani2026deeptime} found that several recent state-of-the-art time-series models
(including TimesNet, iTransformer, and TimeMixer) are often beaten by much simpler
recurrent and convolutional models once tested on real NEM data with actual volatility,
structural breaks and negative prices. This shows that a bigger or more complex model
isn't automatically better; what matters instead is whether a model's assumptions, like
explicit spatial structure, sensible loss design, or multi-task regularisation, fit the
statistical behaviour of the data it's being used on.

\section{Feasibility}

\textbf{Data.} Both reference papers draw on data from AEMO, which publishes historical
Regional Reference Price (RRP), demand, generation, and inter-regional power-flow data
at 5-minute resolution for all five NEM regions, as well as outputs from its ahead-process
forecasting systems. \citet{gani2026deeptime} retrieve RRP data programmatically from
AEMO using the \texttt{nemdatatools} Python library and share their source code, dataset,
and supplementary materials on GitHub, while \citet{li2025mgaal} use regional price,
demand, generation, power-flow, 28-day average-price, and time features. We will use
\texttt{NEMOSIS}, a Python package for downloading and processing historical AEMO market
data, to retrieve the AEMO tables needed to construct realised price, demand, generation,
and inter-regional flow series, and \texttt{NEMSEER} to obtain historical AEMO
ahead-process forecast data such as pre-dispatch outputs. These sources will be combined
and timestamp-aligned to build the feature set for all five regions, giving us a large,
recent, publicly available dataset without writing custom data-download code. \\

\textbf{What has already been done.} Over the past four weeks we have reviewed both
papers in detail, including their data pipelines, architectures, and evaluation
methodology, confirmed that the code and data of \citet{gani2026deeptime} are openly
available as a reproducibility reference, and set up a
\texttt{NEMOSIS}/\texttt{NEMSEER}-based data pipeline as the first step of our own
implementation. \\

\textbf{Planned work.} We have split into parallel workstreams and organised the
implementation into two-week sprints, so each stage exercises concepts covered in
lectures. Sprint~1 (Weeks~6--7, 31/08--13/09) covers data retrieval and preprocessing
with \texttt{NEMOSIS} and \texttt{NEMSEER}; implementation of an MLP, a Temporal
Convolutional Network, a Transformer, an LSTM, and a Gradient Boosting baseline as a
non-deep-learning reference point; a shared evaluation pipeline reporting MAE, RMSE,
sMAPE, and directional accuracy so every model is scored consistently; and finalising
this proposal. This spans course content from data preparation and model selection
through to gradient-based learning, convolutional and recurrent architectures, and
attention mechanisms. Sprints~2--3 (Weeks~8--9) then consolidate evaluation across the
Sprint~1 models and move on to the harder graph-attention multi-task model inspired by
\citet{li2025mgaal}, applying course content on optimisation and regularisation before
building the graph-neural-network layers and the multi-task loss with learned task
weights. \\

\textbf{Evaluation.} We will report MAE, RMSE, sMAPE, and mean directional accuracy on
day-ahead and two-day-ahead horizons, with a dedicated breakdown for the upper/lower 5\%
price tail and negative-price intervals, following \citet{gani2026deeptime}, and
weighted-F1 for the auxiliary spike-classification task once the multi-task model is
implemented, following \citet{li2025mgaal}.\\
\textbf{Risks and mitigation.} The main risks are compute budget for training six model
families across five regions and two horizons, mitigated by training on UQ's Rangpur HPC
cluster for GPU resources well beyond a personal machine, and the added complexity of
the graph-attention model, which is deliberately scheduled for Weeks~8--9 so a working
set of simpler baselines and a shared evaluation pipeline already exist first. Given the
public data, the \texttt{NEMOSIS}/\texttt{NEMSEER} pipeline, the available reference
implementation from \citet{gani2026deeptime}, and the sprint plan's alignment with weekly
course content, we consider the project feasible within the semester timeline.

\smaller
\printbibliography

\end{document}